\input{../doc-class-cours.tex}

\begin{document}

\begin{center}
  \textit{Les propositions mathématiques sont reçues comme vraies parce que personne n’a intérêt qu’elles soient fausses.} - \textbf{Montesquieu}
\end{center}


\subsection*{Ex 1 - Urnes}

On dispose d'une urne A contenant 6 boules numérotées : \quad 7~;~10~;~10~;~12~;~12;~15; \\
et d'une urne B contenant 9 boules numérotées :\quad  2~;~6~;~6~;~6~;~10~;~10~;~24~;~24~;~25. \\
Les boules sont indiscernables au toucher. \\

\medskip

	\begin{enumerate}
		\item On tire une boule dans l'urne A, quelle est la probabilité d'obtenir 12 ?
		\item On tire une boule dans l'urne B, quelle est la probabilité d'obtenir 6 ?
		\item On tire une boule dans l'urne A, quelle est la probabilité d'obtenir un nombre pair ?
		\item On tire une boule dans l'urne B, quelle est la probabilité de ne pas obtenir un nombre pair ?
		\item On tire une boule au hasard dans l'une des urnes. Dans quelle urne la probabilité d'obtenir 10 est la plus grande ?
		\item On verse le contenu de l'urne A dans l'urne B. On tire une boule dans cette urne. Quelle est la probabilité d'obtenir 24 ?
	\end{enumerate}

\subsection*{Ex 2 - Digicode}

Un digicode commande l'ouverture de la porte d'entrée de la maison de la grand-mère de Léna. \\
Léna a oublié le code. Elle sait qu'il est composé d'une lettre A, B, C, ou D suivie d'un chiffre compris entre 1 et 9.\\

\begin{enumerate}
  \item Proposer deux codes différents que Léna peut tester.
  \item Quelle est la probabilité que la grand-mère de Léna ait choisi la lettre C dans son code ?
  \item Montrer que la probabilité que la grand-mère de Léna ait choisi le chiffre 7 dans son code est $\dfrac{1}{9}$.
  \item Léna se souvient que sa grand-mère, enseignante de mathématiques à la retraite, aime bien les nombres premiers. Quelle est la probabilité que le code choisi par sa grand-mère comporte un nombre premier ?
  \item Léna décide de tester tous les codes possibles. Elle estime qu'il lui faut 5 secondes pour essayer un code. Réussira-t-elle à ouvrir la porte de la maison en moins de 4 minutes?
\end{enumerate}

\subsection*{Ex 3 - Musique}

Pour faire écouter de la musique à son enfant, Aurélie a sélectionné 22 chansons : \\
9 chants de Noël, 6 comptines et des berceuses.\\
Le temps d'écoute total des chansons de sa liste est de 55 minutes.

\begin{enumerate}
  \item Calculer le nombre de berceuses présentes dans la liste.
  \item Calculer la durée moyenne d'une chanson de cette liste. Le résultat sera donné en
  minute et seconde.
  \item Aurélie écoute une chanson. Elle utilise la fonction aléatoire de son lecteur, c'est-à-dire que la chanson écoutée est choisie au hasard parmi toutes les chansons de la liste.
	\begin{enumerate}
		\item Quelle est la probabilité que la chanson écoutée soit une comptine ?
		\item Quelle est la probabilité que la chanson écoutée ne soit pas une berceuse ?
		\item Les chansons sont numérotées de 1 à 22. On considère l'évènement :
		
    \begin{center}\og Le numéro de la chanson écoutée est un nombre premier. \fg\end{center}
    La probabilité de cet évènement est-elle supérieure à $\dfrac{1}{3}$ ? Justifier.
  \end{enumerate}
\end{enumerate}

\subsection*{Ex 4 - Dans un cercle}

Sur la figure ci-dessous, on a :

\begin{itemize}
  \item $\mathcal{C}$ est un cercle de centre O et de rayon $4,5$~cm ;
  \item $[\text{AB}]$ est un diamètre de ce cercle et D est un point du cercle ;
  \item les points B, E, A sont alignés, ainsi que les points D, F{}, A ;
  \item les droites (BD) et (EF) sont parallèles ;
  \item BD $= 5,4$ cm ~;~ DA $= 7,2$ cm \:et\: AE $= 2,7$ cm.
  \item le triangle ABD est rectangle en D
\end{itemize}


\begin{figure}[H]
    \centering
    \includegraphics[width=0.5\linewidth]{3x5-proba/ex4.pdf}
\end{figure}


\medskip

\begin{enumerate}
\item Calculer AB{}.
\item Calculer AF{}.
\item 
	\begin{enumerate}
		\item Calculer l’aire du triangle ABD.
		\item Calculer l’aire du disque.
	\end{enumerate}
	
\emph{Rappel} : l’aire du disque est égale à $\pi \times R^2$, où $R$ est le rayon du disque. 
\end{enumerate}

\bigskip

\end{document}